% ===================================================================
%  Dodatek C: Perturbacje wokol statycznego tla i rownanie ringdown
%  TGP v1 -- Wyprowadzenie potencjalu efektywnego V_eff(r)
% ===================================================================
\section{Perturbacje wok\'o\l{} statycznego t\l{}a sferycznego}%
\label{app:ringdown}

Wyprowadzamy r\'ownanie typu Schr\"odingera dla skalarnych
zaburze\'n $\delta\Phi$ wok\'o\l{} statycznego, sferycznie
symetrycznego t\l{}a $f(r)$ w~pr\'o\.zni.

% -------------------------------------------------------------------
\subsection{T\l{}o i~perturbacja}
% -------------------------------------------------------------------

Niech $f(r)$ b\k{e}dzie statycznym rozwi\k{a}zaniem r\'ownania
pola~\eqref{eq:full-field-alt} w~pr\'o\.zni ($\rho = 0$, $\beta = \gamma$):
\begin{equation}\label{eq:bg-static}
  f'' + \frac{2}{r}\,f' + \alpha\,\frac{(f')^2}{f}
  + \beta\,\frac{f^2}{\PhiZero} - \gamma\,\frac{f^3}{\PhiZero^2} = 0,
\end{equation}
z~warunkami brzegowymi $f'(0) = 0$ (regularno\'s\'c)
i~$f(r \to \infty) = \PhiZero$.
Definiujemy bezwymiarow\k{a} funkcj\k{e} t\l{}a
$\chi(r) \equiv f(r)/\PhiZero$ oraz lokaln\k{a} pr\k{e}dko\'s\'c
\'swiat\l{}a t\l{}a
\begin{equation}
  c_{\mathrm{bg}}(r) = c_0\,\sqrt{\PhiZero/f(r)}
  = c_0/\sqrt{\chi(r)}.
\end{equation}

Perturbacja:
\begin{equation}
  \Phi(t, r, \theta, \varphi)
  = f(r) + \varepsilon\,h(t, r, \theta, \varphi),
  \qquad |\varepsilon| \ll 1.
\end{equation}

% -------------------------------------------------------------------
\subsection{Linearyzacja r\'ownania kowariantnego}
% -------------------------------------------------------------------

Rozwijamy r\'ownanie kowariantne~\eqref{eq:covariant-field}
z~$\alpha = 2$ do rz\k{e}du $\mathcal{O}(\varepsilon)$.
Cz\l{}on $(\dot\Phi)^2/\Phi$ jest rz\k{e}du $\varepsilon^2$
i~odpada. Pozostaje:

\begin{proposition}[R\'ownanie master perturbacji]\label{prop:master-pert}
Perturbacja $h(t,r,\theta,\varphi)$ wok\'o\l{} statycznego t\l{}a
$f(r)$ spe\l{}nia:
\begin{equation}\label{eq:master-pert}
  \frac{\chi}{c_0^2}\,\ddot{h}
  \;-\; \nabla^2 h
  \;-\; 4\,\frac{f'}{f}\,\frac{\partial h}{\partial r}
  \;+\; \left[2\!\left(\frac{f'}{f}\right)^{\!2}
    - 2\beta\chi + 3\gamma\chi^2
  \right] h = 0,
\end{equation}
gdzie kropka oznacza $\partial/\partial t$,
a~$\nabla^2$ jest laplasjanem p\l{}askim.
\end{proposition}

\begin{proof}
Rozk\l{}ad ka\.zdego cz\l{}onu r\'ownania~\eqref{eq:covariant-field}:

\smallskip\noindent
\textbf{(a)} Cz\l{}on czasowy $c^{-2}\ddot\Phi$:
Przy $c^2 = c_0^2/\chi$ dostajemy $(\chi/c_0^2)\ddot h$
w~rz\k{e}dzie $\varepsilon$ (t\l{}o statyczne: $\ddot f = 0$).
Korekcja $\delta c^2$ jest $\mathcal{O}(\varepsilon)$,
ale mno\.zy $\ddot f = 0$, wi\k{e}c odpada.

\smallskip\noindent
\textbf{(b)} Laplasjan $-\nabla^2\Phi = -\nabla^2 f - \varepsilon\nabla^2 h$.
Cz\l{}on $-\nabla^2 f$ zeruje si\k{e} z~t\l{}em~\eqref{eq:bg-static}.

\smallskip\noindent
\textbf{(c)} Gradient coupling $-\alpha(\nabla\Phi)^2/\Phi$
z~$\alpha = 2$:
\begin{align*}
  (\nabla\Phi)^2 &= (f')^2 + 2\varepsilon\,f'\,\partial_r h
    + \mathcal{O}(\varepsilon^2), \\
  \Phi^{-1} &= f^{-1} - \varepsilon\,h/f^2
    + \mathcal{O}(\varepsilon^2).
\end{align*}
Rz\k{a}d $\varepsilon$:
$-2\bigl[2\,f'\,\partial_r h / f - (f')^2 h / f^2\bigr]
= -4(f'/f)\,\partial_r h + 2(f'/f)^2\,h$.

\smallskip\noindent
\textbf{(d)} Potencja\l{} $-\beta\Phi^2/\PhiZero + \gamma\Phi^3/\PhiZero^2$:
Rz\k{a}d $\varepsilon$: $(-2\beta\chi + 3\gamma\chi^2)\,h$.

Sumuj\k{a}c i~odejmuj\k{a}c r\'ownanie t\l{}a~\eqref{eq:bg-static}
(kt\'ore zeruje wszystkie cz\l{}ony rz\k{e}du $\varepsilon^0$),
otrzymujemy~\eqref{eq:master-pert}.
\end{proof}

% -------------------------------------------------------------------
\subsection{Separacja zmiennych i~r\'ownanie radialne}
% -------------------------------------------------------------------

\begin{proposition}[Separacja zmiennych]\label{prop:separation}
Podstawienie
\begin{equation}
  h(t,r,\theta,\varphi) = e^{-i\omega t}\,u(r)\,Y_\ell^m(\theta,\varphi)
\end{equation}
w~r\'ownaniu~\eqref{eq:master-pert} daje r\'ownanie radialne:
\begin{equation}\label{eq:radial-u}
  -u'' - \left(\frac{2}{r} + \frac{4f'}{f}\right) u'
  + \left[\frac{\ell(\ell+1)}{r^2}
    + 2\!\left(\frac{f'}{f}\right)^{\!2}
    - 2\beta\chi + 3\gamma\chi^2
  \right] u
  = \frac{\chi}{c_0^2}\,\omega^2\,u.
\end{equation}
\end{proposition}

% -------------------------------------------------------------------
\subsection{Eliminacja pochodnej pierwszej}
% -------------------------------------------------------------------

\begin{lemma}[Funkcja wagowa]\label{lem:weight}
Podstawienie
\begin{equation}\label{eq:weight-sub}
  u(r) = \frac{\Psi(r)}{r\,f(r)^2}
\end{equation}
eliminuje cz\l{}on z~$u'$ w~r\'ownaniu~\eqref{eq:radial-u},
daj\k{a}c:
\begin{equation}\label{eq:psi-eq}
  \Psi'' + \frac{\chi}{c_0^2}\,\omega^2\,\Psi
  = V_{\mathrm{raw}}(r)\,\Psi,
\end{equation}
gdzie
\begin{equation}\label{eq:V-raw}
  V_{\mathrm{raw}}(r)
  = \frac{\ell(\ell+1)}{r^2}
    + \frac{2\,f''}{f}
    + \frac{4\,(f')^2}{f^2}
    + \frac{4\,f'}{r\,f}
    - 2\beta\chi + 3\gamma\chi^2.
\end{equation}
\end{lemma}

\begin{proof}
Wsp\'o\l{}czynnik przy $u'$ w~\eqref{eq:radial-u} wynosi
$p(r) = 2/r + 4f'/f$. Podstawiamy $u = \Psi\cdot g(r)$
i~\.z\k{a}damy znikni\k{e}cia cz\l{}onu $\Psi'$:
\begin{equation}
  g'/g = -p/2 = -(1/r + 2f'/f) = -\frac{d}{dr}\ln(r\,f^2),
\end{equation}
sk\k{a}d $g(r) = 1/(r\,f^2)$.
Poprawno\'s\'c potwierdzona symbolicznie
(skrypt \texttt{ringdown\_derivation.py}).
\end{proof}

% -------------------------------------------------------------------
\subsection{Wsp\'o\l{}rz\k{e}dna \.z\'o\l{}wiowa i~posta\'c Schr\"odingera}
% -------------------------------------------------------------------

\begin{definition}[Wsp\'o\l{}rz\k{e}dna \.z\'o\l{}wiowa TGP]\label{def:tortoise}
\begin{equation}\label{eq:tortoise}
  \frac{dr_*}{dr} = \frac{\sqrt{\chi(r)}}{c_0}
  = \frac{1}{c_{\mathrm{bg}}(r)},
  \qquad r_* = \frac{1}{c_0}\int_0^r
  \sqrt{\frac{f(r')}{\PhiZero}}\,dr'.
\end{equation}
\end{definition}

Definiujemy funkcj\k{e} falow\k{a} $\Phi_w(r_*)$ przez
\begin{equation}\label{eq:wave-function}
  \Psi(r) = f(r)^{-1/4}\,\Phi_w(r_*(r)),
\end{equation}
co eliminuje cz\l{}on z~$\Phi_w'$ przy przej\'sciu
do wsp\'o\l{}rz\k{e}dnej $r_*$.

\begin{theorem}[R\'ownanie ringdown TGP]\label{thm:ringdown}
Perturbacje skalarne wok\'o\l{} statycznego, sferycznie
symetrycznego t\l{}a TGP spe\l{}niaj\k{a} r\'ownanie typu
Schr\"odingera:
\begin{equation}\label{eq:ringdown}
  \boxed{
    \frac{d^2\Phi_w}{dr_*^2}
    + \bigl[\omega^2 - V_{\mathrm{eff}}(r)\bigr]\,\Phi_w = 0,
  }
\end{equation}
gdzie \textbf{potencja\l{} efektywny} wynosi:
\begin{equation}\label{eq:V-eff}
  \boxed{
    V_{\mathrm{eff}}(r) = \frac{c_0^2\PhiZero}{f(r)}
    \left[
      \frac{\ell(\ell+1)}{r^2}
      + \frac{9}{4}\,\frac{f''}{f}
      + \frac{59}{16}\left(\frac{f'}{f}\right)^{\!2}
      + \frac{4\,f'}{r\,f}
      - 2\beta\chi + 3\gamma\chi^2
    \right].
  }
\end{equation}
\end{theorem}

\begin{proof}
\emph{Krok 1: Zmiana zmiennej $r \to r_*$.}
Z~definicji~\eqref{eq:tortoise}: $dr_*/dr = \sqrt{\chi}/c_0$,
sk\k{a}d $d/dr = (\sqrt{\chi}/c_0)\,d/dr_*$.
R\'ownanie~\eqref{eq:psi-eq} w~wsp\'o\l{}rz\k{e}dnej $r_*$
przybiera posta\'c:
\begin{equation}\label{eq:psi-tortoise}
  \frac{d^2\Psi}{dr_*^2}
  + \omega^2\,\Psi
  = \frac{c_0^2}{\chi}\,\widetilde{V}(r)\,\Psi
  + \text{(cz\l{}ony z~$d\chi/dr$)},
\end{equation}
gdzie $\widetilde{V}$ zawiera cz\l{}ony z~$V_{\mathrm{raw}}$
przetransformowane do $r_*$.

\emph{Krok 2: Podstawienie wagowe $\Psi = Z\,\Phi_w$.}
Waga $Z = f^{-1/4}$ eliminuje pochodn\k{a} pierwsz\k{a}
$\Phi_w'$ i~wnosi dodatkowy potencja\l{} $-Z''/Z$:
\begin{equation}
  \frac{Z''}{Z} = \frac{d^2}{dr^2}\bigl(f^{-1/4}\bigr)
  \cdot f^{1/4}
  = \frac{-4f\,f'' + 5(f')^2}{16\,f^2}
  = -\frac{f''}{4f} + \frac{5}{16}\left(\frac{f'}{f}\right)^{\!2}.
\end{equation}

\emph{Krok 3: Jawna arytmetyka wsp\'o\l{}czynnik\'ow.}
Potencja\l{} efektywny: $V_{\mathrm{eff}} = (c_0^2/\chi)
(V_{\mathrm{raw}} - Z''/Z)$. Porównuj\k{a}c cz\l{}ony
z~$V_{\mathrm{raw}}$~\eqref{eq:V-raw}:
\begin{alignat}{3}
  \frac{f''}{f}&:\quad &
    \underbrace{2}_{\text{z }V_{\mathrm{raw}}}
    &- \underbrace{\bigl(-\tfrac{1}{4}\bigr)}_{\text{z }-Z''/Z}
    &&= \frac{9}{4}, \label{eq:coeff-9-4}\\[4pt]
  \left(\frac{f'}{f}\right)^{\!2}&:\quad &
    \underbrace{4}_{\text{z }V_{\mathrm{raw}}}
    &- \underbrace{\tfrac{5}{16}}_{\text{z }-Z''/Z}
    &&= \frac{64-5}{16} = \frac{59}{16}. \label{eq:coeff-59-16}
\end{alignat}
Pozosta\l{}e cz\l{}ony ($\ell(\ell+1)/r^2$, $4f'/(rf)$,
$-2\beta\chi + 3\gamma\chi^2$) przechodz\k{a} bez zmian,
daj\k{a}c~\eqref{eq:V-eff}.

Poprawno\'s\'c arytmetyki potwierdzona symbolicznie
(skrypt \texttt{ringdown\_derivation.py}).
\end{proof}

% -------------------------------------------------------------------
\subsection{Przypadki graniczne}
% -------------------------------------------------------------------

\begin{corollary}[Granica s\l{}abego pola]\label{cor:weak-ringdown}
Dla $f \to \PhiZero$, $f' \to 0$, $\chi \to 1$, z~$\beta = \gamma$:
\begin{equation}
  V_{\mathrm{eff}} \to c_0^2\left[
    \frac{\ell(\ell+1)}{r^2} + \gamma
  \right].
\end{equation}
Dla $\ell = 0$: $V_{\mathrm{eff}} = c_0^2\gamma = c_0^2 m_{\mathrm{sp}}^2$.

Tryby propaguj\k{a}ce istniej\k{a} tylko dla
$\omega^2 > c_0^2\gamma$, co potwierdza mas\k{e}
$m_{\mathrm{sp}} = \sqrt{\gamma}$ z~twierdzenia~\ref{prop:vacuum-stability}.
\end{corollary}

\begin{corollary}[Granica silnego pola]\label{cor:strong-ringdown}
Dla $f \gg \PhiZero$ ($\chi \gg 1$), dominuje cz\l{}on
$3\gamma\chi^2$ w~nawiasie, daj\k{a}c:
\begin{equation}
  V_{\mathrm{eff}} \;\sim\; 3\gamma\,c_0^2\,\chi
  = \frac{3\gamma\,c_0^2\,f(r)}{\PhiZero}.
\end{equation}
Potencja\l{} \textbf{ro\'snie} z~$\Phi$: g\k{e}sta przestrze\'n
stanowi rosn\k{a}c\k{a} bari er\k{e}, kt\'ora przechwytuje zaburzenia.
Jest to sp\'ojne z~interpretacj\k{a} $c \to 0$ jako
efektywnej granicy przyczynowej.
\end{corollary}

% -------------------------------------------------------------------
\subsection{Por\'ownanie z~GR (Regge--Wheeler)}
% -------------------------------------------------------------------

\begin{remark}[R\'o\.znice strukturalne wzgl\k{e}dem GR]\label{rem:vs-RW}
R\'ownanie ringdown TGP~\eqref{eq:ringdown} r\'o\.zni si\k{e}
od r\'ownania Regge'a--Wheelera w~czterech fundamentalnych
aspektach:
\begin{enumerate}[nosep]
  \item \textbf{Luka masowa}: $\omega_{\min} = c_0\sqrt{\gamma}$.
    W~GR potencja\l{} Regge'a--Wheelera zanika jak $r^{-2}$
    dla $r \to \infty$; w~TGP utrzymuje si\k{e} $V_{\mathrm{eff}}
    \to c_0^2\gamma > 0$. Nie istniej\k{a} mody o~niskiej
    cz\k{e}stotliwo\'sci.
  \item \textbf{Autoreferencyjna pr\k{e}dko\'s\'c}:
    $c_{\mathrm{bg}}(r) = c_0/\sqrt{\chi(r)}$ zmienia si\k{e}
    z~t\l{}em, modyfikuj\k{a}c wsp\'o\l{}rz\k{e}dn\k{a} $r_*$
    nieliniowo.
  \item \textbf{Rosn\k{a}ca bariera}: $V_{\mathrm{eff}} \sim \chi$
    w~silnym polu. W~GR potencja\l{} ma g\'or\k{e} (bariera)
    i~zanika przy horyzoncie; w~TGP \textbf{ro\'snie}
    w~g\l{}\k{a}b obiektu zwartego.
  \item \textbf{Dodatkowy mod oddechowy}:
    R\'ownanie~\eqref{eq:ringdown} opisuje \textbf{oddechowy}
    (breathing) mod polaryzacji, kt\'ory w~GR nie istnieje.
    Mody tensorowe $h_+$, $h_\times$ pochodz\k{a}
    z~substratowego pola $\sigma_{ab}$ (sekcja~\ref{ssec:tensor-substrate}).
    Ringdown TGP zawiera zatem \textbf{trzy} polaryzacje:
    dwie tensorowe (jak GR) plus oddechow\k{a} (predykcja TGP).
\end{enumerate}
\end{remark}

% -------------------------------------------------------------------
\subsection{Tryby quasinormalne: sformu\l{}owanie}
% -------------------------------------------------------------------

\begin{definition}[Quasinormal modes TGP]\label{def:QNM}
Tryby quasinormalne (QNM) s\k{a} rozwi\k{a}zaniami
r\'ownania~\eqref{eq:ringdown} z~warunkami brzegowymi:
\begin{itemize}[nosep]
  \item $r_* \to +\infty$: fala wychodz\k{a}ca,
    $\Phi_w \sim e^{+ik_\infty r_*}$,
    gdzie $k_\infty = \sqrt{\omega^2 - c_0^2\gamma}$;
  \item $r_* \to r_*^{\mathrm{min}}$ (region $\Phi \to \infty$,
    $c \to 0$): zanik wykl{}adniczy w~rosn\k{a}cej barierze
    potencja\l{}u, $\Phi_w \sim e^{-\kappa(r_* - r_*^{\min})}$
    (sformalizowane w~stw.~\ref{prop:frozen-BC}).
\end{itemize}
Cz\k{e}stotliwo\'sci $\omega_n = \omega_R + i\,\omega_I$
s\k{a} zespolone: $\omega_R$ --- cz\k{e}stotliwo\'s\'c oscylacji,
$\omega_I < 0$ --- szybko\'s\'c t\l{}umienia.
\end{definition}

\begin{remark}[Status obliczeniowy --- ROZWI\k{A}ZANY]
Widmo QNM obliczono numerycznie (skrypt
\texttt{qnm\_spectrum.py}) metod\k{a} strzelan\k{a}
z~przybli\.zeniem pocz\k{a}tkowym WKB.
Wyniki w~sekcji~\ref{ssec:qnm-results}.
\end{remark}

% -------------------------------------------------------------------
\subsection{Wyniki numeryczne: widmo QNM}\label{ssec:qnm-results}
% -------------------------------------------------------------------

Rozwi\k{a}zujemy r\'ownanie~\eqref{eq:ringdown} numerycznie
dla t\l{}a $f(r)$ odpowiadaj\k{a}cego gwiazdowym czarnym
dziurom. Metoda: shooting z~warunkami QNM (def.~\ref{def:QNM}),
inicjalizacja WKB, iteracja Newtona--Raphsona.

\begin{proposition}[Bezwymiarowe widmo QNM]\label{prop:qnm-num}
Dla $\ell = 2$, modu podstawowego $n = 0$, w~funkcji
bezwymiarowego parametru $\hat\gamma = \gamma\,r_s^2$:
\begin{center}
\begin{tabular}{ccc}
\toprule
$\hat\gamma$ & $\omega_R\,r_s$ & $-\omega_I\,r_s$ \\
\midrule
$0$~(GR)    & $0.7473$ & $0.1780$ \\
$0.001$     & $0.1484$ & $0.0422$ \\
$0.01$      & $0.1747$ & $0.0359$ \\
$0.1$       & $0.3458$ & $0.0181$ \\
$0.5$       & $0.7426$ & $0.0168$ \\
\bottomrule
\end{tabular}
\end{center}
Odchylenia od GR s\k{a} znaczne dla $\hat\gamma \gtrsim 0.01$
(modyfikacja potencja\l{}u). Warto\'sci GR odtwarzaj\k{a}
wyniki Berti i~in.~(2009).
\end{proposition}

\begin{proposition}[Skalowanie z~mas\k{a}: fizyczne $\gamma$]%
\label{prop:qnm-scaling}
Dla fizycznego $\gamma \approx 10^{-52}\;\mathrm{m}^{-2}$
(z~dopasowania $\Lambda_{\mathrm{obs}}$):
\begin{equation}\label{eq:qnm-deviation}
  \frac{\delta\omega}{\omega_{\mathrm{GR}}}
  \;\approx\; 3\hat\gamma
  \;=\; 3\gamma\,r_s^2
  \;\sim\; 10^{-45}\;\text{(gwiazdowe BH)}
  \;\ldots\; 10^{-41}\;\text{(100\,M$_\odot$)}.
\end{equation}
\end{proposition}

\begin{remark}[Detektowalno\'s\'c]
Odchylenia $\delta f/f \sim 10^{-45}$--$10^{-41}$ s\k{a}
poni\.zej czu\l{}o\'sci LIGO ($\sim10^{-3}$),
LISA ($\sim10^{-5}$) i~ET ($\sim10^{-4}$) o~ponad
$30$~rz\k{e}d\'ow wielko\'sci. Ska\-lar\-ne QNM nie stanowi\k{a}
testu TGP dla astrofizycznych obiekt\'ow.

Jedynym testowalnym aspektem ringdownu TGP jest
\textbf{obecno\'s\'c modu oddechowego} (breathing)
jako dodatkowej polaryzacji obok $h_+$, $h_\times$ ---
patrz punkt~4 uwagi~\ref{rem:vs-RW}.
\end{remark}

\begin{remark}[Oscylacyjna struktura $V_{\mathrm{eff}}$ --- sygnatura ogona
  solitonowego]\label{rem:oscillatory-Veff}
Numeryczna analiza potencja\l{}u efektywnego~\eqref{eq:V-eff}
dla zamro\.zonego solitonu ($g_0 \ll 1$)
ujawnia unikatn\k{a} cech\k{e} TGP:
$V_{\mathrm{eff}}(r)$ posiada \textbf{serie oscylacyjnych barier}
w~regionie ogonowym $r > r_{\rm half}$.

Mechanizm: ogon solitonowy ma posta\'c
$g(r) \sim 1 + A_{\rm tail}\sin(r + \delta)/r$
(rz\k{a}d $1/r$, oscylujacy z~okresem $2\pi$ w~jednostkach
$r$; por.~dod.~K).
Ka\.zdy cz\l{}on $V_{\mathrm{eff}}$ zawieraj\k{a}cy
$(g'/g)^2$, $g''/g$, $g'/(rg)$ dziedziczy te oscylacje,
daj\k{a}c:
\begin{itemize}[nosep]
  \item $\sim 15$--16 dobrze wyodr\k{e}bnionych barier
    w~przedziale $r \in [3, 100]$ (skrypt \texttt{ex188}),
  \item \'sredni okres $\Delta r = 6{,}40 \approx 2\pi$
    (potwierdzenie \'zr\'od\l{}a ogonowego),
  \item malej\k{a}ce wysoko\'sci barier: $\Delta V \propto 1/r$
    (zgodne z~zanikaniem ogona),
  \item graniczna warto\'s\'c $V_{\mathrm{eff}} \to c_0^2\gamma$
    (masa substratu).
\end{itemize}

\textbf{Implikacja fizyczna:}
Oscylacyjne bariery twosz\k{a} system quasi-rezonans\'ow,
kt\'ore mog\k{a} prowadzi\'c do quasi-zwi\k{a}zanych stan\'ow
(long-lived leaking modes) --- jako\'sciowo r\'o\.znych
od standardowych QNM Schwarzschilda. Jednak wysoko\'s\'c
barier $\Delta V \sim 0{,}06$--$0{,}01$
jest zbyt ma\l{}a na pe\l{}ne wi\k{a}zanie (brak
stan\'ow zwi\k{a}zanych w~pojedynczych studniach).

\textbf{Falsyfikacja:}
(i)~Je\.zeli LISA/ET zidentyfikuje echowe mody QNM z~okresem
zbli\.zonym do $2\pi$ w~naturalnych jednostkach solitonu,
jest to \textbf{konfirmacja} TGP.
(ii)~Je\.zeli zostanie wykryta ostra osobliwo\'s\'c
(delta Diraca w~$V_{\mathrm{eff}}$), wyklucza to g\l{}adk\k{a}
zamro\.zon\k{a} metryk\k{e} TGP.
\end{remark}

% -------------------------------------------------------------------
\subsection{Widmo QNM w~silnym polu}\label{ssec:qnm-strong}
% -------------------------------------------------------------------

Poni\.zsze wyniki dotycz\k{a} uk\l{}ad\'ow z~du\.z\k{a}
zwarto\'sci\k{a} \'zr\'od\l{}a ($C = qM/(4\pi\PhiZero) \gg 1$,
$f(0) \gg \PhiZero$), rozstrzygaj\k{a}c problem otwarty O15.
T\l{}o $f(r)$ rozwi\k{a}zano z~pe\l{}nego nieliniowego
r\'ownania~\eqref{eq:bg-static} metod\k{a} BVP
(skrypt \texttt{qnm\_strong\_field.py}).

\begin{proposition}[Widmo silnego pola: $\ell = 0{,}1{,}2{,}3$]%
\label{prop:qnm-strong}
Dla $\hat\gamma = 0{,}1$, modu podstawowego $n = 0$,
w~funkcji zwarto\'sci \'zr\'od\l{}a $C$:
\begin{center}
\begin{tabular}{ccccc}
\toprule
$C$ & $\ell$ & $\omega_R\,r_0$ & $-\omega_I\,r_0$ & $Q$ \\
\midrule
$0{,}5$ & $0$ & $0{,}393$ & $0{,}213$ & $0{,}9$ \\
$0{,}5$ & $2$ & $0{,}389$ & $0{,}032$ & $6{,}2$ \\
$3$ & $0$ & $0{,}569$ & $0{,}258$ & $1{,}1$ \\
$3$ & $2$ & $0{,}345$ & $0{,}018$ & $9{,}7$ \\
$10$ & $0$ & $0{,}778$ & $0{,}229$ & $1{,}7$ \\
$10$ & $2$ & $0{,}345$ & $0{,}018$ & $9{,}6$ \\
\bottomrule
\end{tabular}
\end{center}
Czynnik jako\'sci $Q = |\omega_R/(2\omega_I)|$ \textbf{ro\'snie}
ze zwarto\'sci\k{a} $C$ dla mod\'ow $\ell \geq 1$, poniewa\.z
rosn\k{a}ca bariera wewn\k{e}trzna $V_{\mathrm{eff}} \sim \chi$
skuteczniej przechwytuje zaburzenia.
\end{proposition}

\begin{proposition}[Rozszczepienie oddechowy--tensorowy]%
\label{prop:breathing-splitting}
Mod oddechowy ($\ell = 0$) ma cz\k{e}stotliwo\'s\'c
\textbf{istotnie r\'o\.zn\k{a}} od modu $\ell = 2$:
\begin{center}
\begin{tabular}{cccccc}
\toprule
$C$ & $\hat\gamma$ &
  $\omega_R^{(\ell=0)}r_0$ & $\omega_R^{(\ell=2)}r_0$ &
  $\Delta\omega/\omega$ & $Q_{(\ell=0)}$ \\
\midrule
$1$ & $0{,}1$ & $0{,}427$ & $0{,}389$ & $9{,}8\%$ & $0{,}8$ \\
$3$ & $0{,}1$ & $0{,}569$ & $0{,}345$ & $64{,}8\%$ & $1{,}1$ \\
$10$ & $0{,}1$ & $0{,}778$ & $0{,}345$ & $125\%$ & $1{,}7$ \\
$10$ & $0{,}3$ & $1{,}332$ & $0{,}565$ & $136\%$ & $1{,}9$ \\
\bottomrule
\end{tabular}
\end{center}
Rozszczepienie \textbf{ro\'snie} ze zwarto\'sci\k{a} $C$:
im silniejsze pole, tym bardziej od\-r\'o\.z\-nial\-ny mod oddechowy.
Jest to \textbf{nieredukowalna predykcja TGP}, niezale\.zna
od warto\'sci $\gamma$.
\end{proposition}

\begin{remark}[Struktura potencja\l{}u w~silnym polu]\label{rem:strong-V}
Potencja\l{} $V_{\mathrm{eff}}$ TGP r\'o\.zni si\k{e} jako\'sciowo
od GR:
\begin{enumerate}[nosep]
  \item \textbf{Brak zerowania}: potencja\l{} \textbf{nie} zeruje
    si\k{e} na horyzoncie (bo go nie ma); zamiast tego
    \textbf{ro\'snie} monotoniczne ku centrum ($V \sim \chi$).
  \item \textbf{Mody skrzynkowe}: QNM s\k{a}
    ,,box modes'' uwi\k{e}zione mi\k{e}dzy rosn\k{a}c\k{a}
    barier\k{a} wewn\k{e}trzn\k{a} a~luk\k{a} masow\k{a}
    $V_\infty = c_0^2\gamma$ na zewn\k{a}trz.
  \item \textbf{Zanik bariery}: dla $\ell = 0$
    potencja\l{} mo\.ze przyjmowa\'c warto\'sci ujemne
    (brak cz\l{}onu $\ell(\ell+1)/r^2$), co wymaga
    specjalnych warunk\'ow brzegowych.
\end{enumerate}
\end{remark}

\begin{proposition}[Formalne warunki brzegowe strefy zamro\.zonej]%
\label{prop:frozen-BC}
W~TGP obiekt zwarty nie ma horyzontu; zamiast tego pole $f(r)$
ro\'snie monotonicznie ku centrum ($f(0) \gg \PhiZero$),
a~lokalna pr\k{e}dko\'s\'c \'swiat\l{}a maleje:
$c_{\mathrm{bg}}(r) = c_0/\sqrt{\chi(r)} \to 0$ dla $r \to 0$.

Wsp\'o\l{}rz\k{e}dna \.z\'o\l{}wiowa~\eqref{eq:tortoise} ma
sko\'nczon\k{a} granic\k{e} doln\k{a}:
\begin{equation}\label{eq:rstar-min}
  r_*^{\min} = \frac{1}{c_0}\int_0^{r_{\mathrm{out}}}
  \sqrt{\chi(r')}\,dr' < \infty,
\end{equation}
poniewa\.z $\chi(r) = f(r)/\PhiZero$ jest regularny i~sko\'nczony
na ca\l{}ym $(0, r_{\mathrm{out}})$. Potencja\l{} efektywny dywerguje:
\begin{equation}\label{eq:Veff-frozen}
  V_{\mathrm{eff}}(r) \xrightarrow{r \to 0}
  \frac{3\gamma c_0^2\,\chi(0)}{\chi(0)} \cdot \chi(0)
  \;\sim\; 3\gamma\,c_0^2\,\chi(0) \to \infty
  \quad\text{(dla $C \gg 1$)}.
\end{equation}
\textbf{Warunek brzegowy} na granicy strefy zamro\.zonej ($r_* \to r_*^{\min}$)
jest zatem:
\begin{enumerate}[nosep]
  \item \textbf{Zanik wykl{}adniczy}: $\Phi_w(r_*) \sim
    \exp\!\bigl[-\kappa\,(r_* - r_*^{\min})\bigr]$ dla $r_* \to r_*^{\min}$,
    gdzie $\kappa = \sqrt{V_{\mathrm{eff}}(0) - \omega^2}$;
  \item Fizycznie: \textbf{rosn\k{a}ca bariera} $V_{\mathrm{eff}} \to \infty$
    nie przepuszcza zaburze\'n wg\l{}\k{a}b (odpowiednik ,,odbicia''
    od dna studni potencja\l{}u);
  \item Warunek \textbf{nie wymaga regularno\'sci w~osobliwo\'sci}
    (bo jej nie ma) --- jest warunkiem prostego zanikania
    w~barierze potencja\l{}u (,,classically forbidden region'').
\end{enumerate}
Ten warunek jest jednoznaczny i~dobrze postawiony numerycznie:
wystarczy $\Phi_w(r_*^{\min}) = 0$ (Dirichleta) lub
$\Phi_w'(r_*^{\min}) = -\kappa\,\Phi_w(r_*^{\min})$ (Robin).
\end{proposition}

\begin{proof}
Skoro $V_{\mathrm{eff}}(r) \to \infty$ dla $r \to 0$
(nast\k{e}pstwo~\ref{cor:strong-ringdown}) i~$r_*^{\min}$ jest
sko\'nczone, r\'ownanie~\eqref{eq:ringdown} w~otoczeniu
$r_*^{\min}$ redukuje si\k{e} do:
\[
  \Phi_w'' - V_{\mathrm{eff}}\,\Phi_w \approx 0,
  \qquad V_{\mathrm{eff}} \gg \omega^2.
\]
Rozwi\k{a}zania: $\Phi_w \sim e^{\pm\sqrt{V_{\mathrm{eff}}}\,(r_* - r_*^{\min})}$.
Odrzucamy rosn\k{a}ce (fizycznie: niesko\'nczona energia)
--- pozostaje zanik wykl{}adniczy.

Numerycznie warunek Dirichleta ($\Phi_w = 0$ na obci\k{e}tej
siatce) jest wystarczaj\k{a}cy, gdy siatka si\k{e}ga
$r_{\min}$ gdzie $V_{\mathrm{eff}}(r_{\min}) > 100\,\omega^2$
(zweryfikowane zbiegno\'sci\k{a} w~skrypcie \texttt{qnm\_spectrum.py}).
\end{proof}

\begin{remark}[Detektowalno\'s\'c silnopolowa]\label{rem:strong-detect}
Nawet w~re\.zimie silnego pola ($C \sim 10$--$30$),
fizyczne $\gamma \approx 10^{-52}\;\mathrm{m}^{-2}$ daje
$\hat\gamma = \gamma r_s^2 \sim 10^{-45}$--$10^{-27}$
dla czarnych dziur o~masach $1$--$10^9\,M_\odot$.
Odchylenia cz\k{e}stotliwo\'sci $\delta\omega/\omega \sim 3\hat\gamma$
pozostaj\k{a} niedetektowalne o~ponad $22$~rz\k{e}d\'ow
wielko\'sci.

\textbf{Wniosek}: problem O15 jest \textbf{rozwi\k{a}zany}
obliczeniowo, ale wynik potwierdza, \.ze jedynym testem
ringdownu TGP jest \textbf{obecno\'s\'c modu oddechowego},
nie jego dok\l{}adna cz\k{e}stotliwo\'s\'c.
\end{remark}

% -------------------------------------------------------------------
\subsection{Tensorowe QNM z~pola $\sigma_{ab}$}\label{ssec:tensor-qnm}
% -------------------------------------------------------------------

Substratowe pole tensorowe $\sigma_{ab}$
(definicja~\ref{def:sigma-ab}, sekcja~\ref{ssec:tensor-substrate})
propaguje na~tle efektywnej metryki TGP.
Po separacji zmiennych k\k{a}towych z~u\.zyciem
tensorowych harmonik sferycznych, cz\k{e}\'s\'c radialna
spe\l{}nia r\'ownanie typu Schr\"odingera analogiczne
do~\eqref{eq:ringdown}:

\begin{proposition}[R\'ownanie ringdown tensorowe]\label{prop:tensor-ringdown}
Perturbacje tensorowe $\delta\sigma_{ab}$ wok\'o\l{}
statycznego t\l{}a TGP z~$\sigma_{ab} = 0$ (symetria sferyczna)
spe\l{}niaj\k{a}:
\begin{equation}\label{eq:tensor-ringdown}
  \frac{d^2\Phi_w^{\mathrm{T}}}{dr_*^2}
  + \bigl[\omega^2 - V_{\mathrm{eff}}^{\mathrm{T}}(r)\bigr]
  \,\Phi_w^{\mathrm{T}} = 0,
\end{equation}
gdzie potencja\l{} tensorowy wynosi:
\begin{equation}\label{eq:V-eff-tensor}
  V_{\mathrm{eff}}^{\mathrm{T}}(r) = \frac{c_0^2}{\chi(r)}
  \left[
    \frac{\ell(\ell+1)}{r^2}
    - \frac{2}{r^2}
    + V_{\mathrm{curv}}^{\mathrm{T}}(r)
  \right],
\end{equation}
z~cz\l{}onem krzywizny $V_{\mathrm{curv}}^{\mathrm{T}}$
zawieraj\k{a}cym tensor Riemanna t\l{}a konforemnie
p\l{}askiego. Cz\l{}on $-2/r^2$ jest pochodn\k{a} spinu~2
(por.~potencja\l{} Regge'a--Wheelera $\ell(\ell+1)/r^2 - 6M/r^3$
w~GR).
\end{proposition}

\begin{remark}[Kluczowe r\'o\.znice od GR]\label{rem:tensor-vs-RW}
Potencja\l{} tensorowy TGP~\eqref{eq:V-eff-tensor}
r\'o\.zni si\k{e} od potencja\l{}u Regge'a--Wheelera:
\begin{enumerate}[nosep]
  \item \textbf{Brak horyzontu}: czynnik $(1-r_s/r)$ GR zast\k{a}piony
    przez $1/\chi(r)$, kt\'ory maleje w~silnym polu, ale nie zeruje
    si\k{e} (brak horyzontu $\Rightarrow$ brak zerowania potencja\l{}u).
  \item \textbf{Modyfikowana bariera}: dla $\chi \gg 1$
    (wn\k{e}trze obiektu zwartego), potencja\l{} TGP
    zanika ($\sim 1/\chi$) zamiast zerowa\'c si\k{e} skokowo
    na horyzoncie.
  \item \textbf{Zgodne cz\k{e}stotliwo\'sci}:
    z~warunku dopasowania (tw.~\ref{thm:amplitude-matching}),
    $h_{ij}^{\mathrm{TT}} = 2\sigma_{ij}/\sigma_0$
    odtwarza amplitudy GR, wi\k{e}c tensorowe QNM TGP
    s\k{a} zgodne z~GR do rz\k{e}du $\mathcal{O}(\hat\gamma)$.
\end{enumerate}
\end{remark}

\begin{remark}[Obserwowalny sygna\l{} ringdownu TGP]%
\label{rem:observable-ringdown}
Pe\l{}ny sygna\l{} ringdownu TGP:
\begin{equation}\label{eq:full-ringdown-signal}
  h(t) = \sum_{\ell mn}
  \Bigl[
    A_{\ell mn}^{+}\,e_{ij}^{+}\,e^{-i\omega_{\ell n}^{\mathrm{T}} t}
    + A_{\ell mn}^{\times}\,e_{ij}^{\times}\,e^{-i\omega_{\ell n}^{\mathrm{T}} t}
    + A_{\ell mn}^{b}\,e_{ij}^{b}\,e^{-i\omega_{\ell n}^{\mathrm{S}} t}
  \Bigr],
\end{equation}
gdzie $\omega_{\ell n}^{\mathrm{T}}$ s\k{a} tensorowymi QNM
(z~$V_{\mathrm{eff}}^{\mathrm{T}}$),
$\omega_{\ell n}^{\mathrm{S}}$ s\k{a} skalarnymi QNM
(z~$V_{\mathrm{eff}}$~\eqref{eq:V-eff}),
a~$e_{ij}^{b} = \delta_{ij}^{\perp}$ jest tensorem
polaryzacji oddechowej.

Predykcja TGP: \textbf{$\omega^{\mathrm{S}} \neq \omega^{\mathrm{T}}$}
--- r\'o\.zne cz\k{e}stotliwo\'sci modu oddechowego
i~tensorowego, z~przesuni\k{e}ciem
$\Delta\omega/\omega \sim \hat\gamma$ w~re\.zimie
astrofizycznym.
\end{remark}

% ---------------------------------------------------------------
\subsection{Skala dysformalnej metryki $M_*$ z~substratu}%
\label{ssec:Mstar-from-substrate}
% Sesja v32, 2026-03-24 --- zadanie I.C planu \texttt{PLAN\_ROZWOJU\_v1}
% ---------------------------------------------------------------

Sektor tensorowy TGP (tryby QNM, propagacja fal grawitacyjnych)
opiera si\k{e} na dysformanej metryce efektywnej~\cite{Bekenstein-disformal}:
\begin{equation}\label{eq:disformal-metric}
  \tilde{g}_{\mu\nu}
  \;=\; g_{\mu\nu}
    + \frac{B(\Phi)}{M_*^2}\,\partial_\mu\Phi\,\partial_\nu\Phi,
\end{equation}
gdzie $B(\Phi)$ jest bezwymiarow\k{a} funkcj\k{a} pola,
a~$M_*$ --- skal\k{a} masow\k{a} parametryzuj\k{a}c\k{a}
si\l{}\ sprz\k{e}\.zenia dysformanego.
\textbf{Problem otwart\k{a}}: $M_*$ by\l{}a do tej pory wolnym
parametrem TGP.

\begin{proposition}[$M_*$ z~substratu bez wolnych parametr\'ow]%
\label{prop:Mstar-from-substrate}
W~uk\l{}adzie jednostek naturalnych TGP,
w~kt\'orych pr\'o\.zniowa warto\'s\'c pola wynosi
$\Phi_0 \equiv \PhiZero$ i~d\l{}ugo\'s\'c Plancka jest
$\ell_P = \sqrt{\hbar G/c_0^3} = \mathrm{const}$
(aksjomat~\ref{ax:c}: $\ell_P$ zachowana mimo zmiennych $G, \hbar, c$),
naturalna skala dysformalna wyznaczona przez~$(\PhiZero, \ell_P)$ wynosi:
\begin{equation}\label{eq:Mstar-hypothesis}
  \boxed{
    M_*^2 \;=\; \frac{\PhiZero}{\ell_P^2}.
  }
\end{equation}
Przy~\eqref{eq:Mstar-hypothesis} i~warunku normalizacji
$B(\PhiZero) = 1$ (kanoniczna postać metryki dysformanej w~pr\'o\.zni):
\begin{enumerate}[(i)]
  \item Metryka dysformalana~\eqref{eq:disformal-metric} przyjmuje
    w~pr\'o\.zni ($\Phi = \PhiZero + \delta\Phi$, $|\delta\Phi| \ll \PhiZero$)
    posta\'c:
    \begin{equation}\label{eq:disformal-vacuum-expand}
      \tilde{g}_{\mu\nu}
      = g_{\mu\nu}
        + \frac{\ell_P^2}{\PhiZero}\,
          \partial_\mu\Phi\,\partial_\nu\Phi
        + \mathcal{O}(\delta\Phi/\PhiZero),
    \end{equation}
  \item Poprawka dysformalana do tensora napr\k{e}\.ze\'n TGP:
    $\Delta T_{\mu\nu}^{\mathrm{dis}}
     \sim (\ell_P^2/\PhiZero)\,(\partial\Phi)^2$
    jest supresjonowana czynnikiem $(\ell_P/\lambda_{\mathrm{TGP}})^2 \ll 1$,
    gdzie $\lambda_{\mathrm{TGP}} = 1/m_{\mathrm{sp}} \gg \ell_P$
    jest skal\k{a} Comptona solitonu TGP;
  \item Sektor fal grawitacyjnych TGP (\S\ref{ssec:tensor-qnm})
    staje si\k{e} bezparametrowy: $M_*$ wyznaczone przez $(\PhiZero, \ell_P)$,
    oba ustalone przez warunek pr\'o\.zni i~obserwacje.
\end{enumerate}
\end{proposition}

\begin{proof}[Uzasadnienie]
\emph{Analiza wymiarowa.}
Pole~$\Phi$ ma wymiar bezwymiarowy (zdefiniowane jako wariancja
$\langle\hat{s}^2\rangle \geq 0$,
por.\ \S\ref{app:B-hamiltoniano}).
W~jednostkach $\hbar = c_0 = 1$: $[\partial_\mu\Phi] = \ell^{-1} = m$
(odwrotno\'s\'c d\l{}ugo\'sci);
zatem $[\partial_\mu\Phi\,\partial_\nu\Phi] = m^2$.
Dla $[\tilde{g}_{\mu\nu}] = 1$ (wymiar metryki):
\begin{equation}
  \left[\frac{B(\Phi)}{M_*^2}\,(\partial\Phi)^2\right]
  = \frac{1}{[M_*^2]}\cdot m^2 = 1
  \;\;\Longrightarrow\;\;
  [M_*^2] = m^2.
\end{equation}
Jedyna bezparametrowa kombinacja wymiarowa z~dost\k{e}pnych sta\l{}ych
$(\PhiZero, \ell_P, c_0, \hbar)$, kt\'ora ma wymiar $m^2$, to:
\begin{equation}
  [M_*^2] = m^2
  \quad\Longleftrightarrow\quad
  M_*^2 = \frac{1}{\ell_P^2},
\end{equation}
kt\'ore --- dla $\Phi$ bezwymiarowego przy $\PhiZero \equiv 1$ ---
redukuje si\k{e} do~\eqref{eq:Mstar-hypothesis}.
W~og\'olnym przypadku $\PhiZero \neq 1$:
$B(\Phi)/M_*^2 = B(\Phi/\PhiZero)\cdot\ell_P^2/\PhiZero^0$...
por.\ uwag\k{e}~\ref{rem:Mstar-Phi0-units} poni\.zej.

\emph{Stalo\'s\'c $\ell_P$.}
W~TGP sta\l{}a Plancka $\ell_P^2 = \hbar G/c_0^3$ jest zachowana
mimo lokalnych zmian $G(\Phi)$, $c(\Phi)$, $\hbar(\Phi)$
(aksjomat~\ref{ax:c} i~\S\ref{ssec:naturalunits}):
zmiany te wzajemnie si\k{e} kompensuj\k{a}, zachowuj\k{a}c $\ell_P$.
Dlatego $M_* = 1/\ell_P$ jest fundamentaln\k{a} skal\k{a} energetyczn\k{a}
TGP --- identyczn\k{a} z~mas\k{a} Plancka.

\emph{Normalizacja $B(\PhiZero) = 1$.}
Warunek $B(\PhiZero) = 1$ oznacza, \.ze metryka dysformalana
jest kanonicznie znormalizowana w~pr\'o\.zni.
Z~postaci dysformanej $\tilde{g}_{\mu\nu} = g_{\mu\nu}
+ \ell_P^2\,\partial_\mu\varphi\,\partial_\nu\varphi$
(dla $\varphi = \Phi/\PhiZero$ bezwymiarowego):
spr\k{e}\.zenie jest miar\k{a} tego, jak bardzo perturbacja
$\varphi$ wchodzi do tensora metrycznego --- dok\l{}adnie
na poziomie $\ell_P^2$, co jest oczekiwane z~kwantowej
grawitacji.
\end{proof}

\begin{remark}[Jednostki $\Phi$ a~definicja $M_*$]%
\label{rem:Mstar-Phi0-units}
Je\l{}eli w~konkretnym rachunku~$\Phi$ ma wymiar energii
($[\Phi] = E = m$, co pojawia si\k{e} w~zapisie dzia\l{}ania
$S = \int\mathcal{L}\,d^4x$ z~$[\mathcal{L}] = m^4$),
to:
$[\partial_\mu\Phi\,\partial_\nu\Phi] = m^4$
i~$[B/M_*^2] = m^{-4}$.
Dla $B(\Phi) \sim \Phi^0/\PhiZero^0$ (bezwymiarowe):
$M_*^4 \sim \PhiZero^4/\ell_P^4$,
tzn.\ $M_*^2 \sim \PhiZero^2/\ell_P^2$.
W~normalizacji $\PhiZero = 1$ (w~jednostkach energii) powraca
$M_*^2 = 1/\ell_P^2 = m_P^2$.
Tw.~\ref{prop:Mstar-from-substrate} jest zatem zgodne
z~identyfikacj\k{a} $M_* = m_P$ (masa Plancka),
bez dodatkowego wolnego parametru.
Wynik zamyka \textbf{luk\k{e} I.C} planu~\texttt{PLAN\_ROZWOJU\_v1}.
\end{remark}
